% mazzi: 06
\chapter{Organizzazione del sistema nervoso}

% ---------------------------------------------------------------------------
\section{La funzione integrativa del sistema nervoso}

Il sistema nervoso deve elaborare nello stesso istante un'enorme quantità di segnali
provenienti dall'ambiente esterno, integrarli con quanto avviene all'interno del corpo e
produrre una risposta congruente agli stimoli: svolge cioè una \textbf{funzione integrativa}.

\rubrica{Esempio: afferrare una pallina in arrivo}
Per compiere questo gesto apparentemente semplice il sistema nervoso deve, in un secondo:
\begin{enumerate}
  \item calcolare \textbf{peso} e \textbf{dimensioni} della pallina;
  \item calcolare la \textbf{velocità} con cui si sta avvicinando;
  \item calcolare la \textbf{posizione} di tutto il corpo (qui schematizzata nella posizione
    della mano che dovrà afferrarla);
  \item prevedere eventuali \textbf{accelerazioni o deviazioni} durante il tragitto
    (p.\,es.\ il vento, se si è all'aperto);
  \item recuperare dal sistema nervoso centrale lo \textbf{schema} (o programma)
    \textbf{motorio} della prensione;
  \item già al momento del contatto pallina-mano, cominciare a \textbf{stringere};
  \item provvedere all'\textbf{alternanza delle unità motorie}: i muscoli mantengono la
    contrazione solo per alcuni millisecondi, quindi altre unità motorie devono subentrare a
    mantenere la posizione quando le prime iniziano a rilassarsi.
\end{enumerate}

Al compito concorrono più aree e sistemi contemporaneamente: \textbf{corteccia visiva},
\textbf{corteccia somatosensoriale}, \textbf{corteccia motoria} e \textbf{sistemi
motivazionali}. Tutto ciò è possibile grazie alla complessità dell'organizzazione del sistema
nervoso.

% ---------------------------------------------------------------------------
\section{Organizzazione generale del sistema nervoso}

\begin{itemize}
  \item \textbf{sistema nervoso centrale} (SNC): encefalo e midollo spinale; riceve l'\textit{input}
    e genera l'\textit{output};
  \item \textbf{sistema nervoso periferico} (SNP): fibre \textbf{afferenti}, che provengono dalla
    periferia del corpo ed entrano nel SNC, e fibre \textbf{efferenti}, che lasciano il SNC per
    dirigersi verso la periferia.
\end{itemize}

Una seconda suddivisione, trasversale alla prima, distingue:
\begin{itemize}
  \item \textbf{sistema nervoso somatico}: fibre afferenti ed efferenti (più il SNC) che si
    occupano del \textit{soma}, cioè di tutto il corpo \textbf{esclusi gli organi interni};
  \item \textbf{sistema nervoso autonomo} (o viscerale): fibre afferenti ed efferenti degli organi
    su cui \textbf{non agisce la volontà}.
\end{itemize}

\rubrica{Vie afferenti (input)}
\begin{itemize}
  \item \textbf{sensibilità somatica}: p.\,es.\ tatto, pressione, dolore;
  \item \textbf{organi di senso speciale}: p.\,es.\ vista e udito;
  \item \textbf{sensibilità viscerale}: da tutti gli organi interni (cuore, polmone, sistema
    gastroenterico, \dots).
\end{itemize}

\rubrica{Vie efferenti (output)}
\begin{itemize}
  \item \textbf{somatiche}: hanno un unico effettore, il \textbf{muscolo scheletrico}, raggiunto
    dalle vie motorie che escono dal midollo spinale;
  \item \textbf{autonome}: raggiungono muscolo cardiaco, muscolo liscio e ghiandole, oltre al
    \textbf{sistema nervoso enterico}, in rapporto reciproco con il tratto gastrointestinale.
\end{itemize}

\rubrica{Le due branche dell'autonomo}
Il sistema nervoso autonomo si suddivide in \textbf{ortosimpatico} (detto semplicemente
\textit{simpatico}) e \textbf{parasimpatico}; tutti gli organi interni ricevono innervazione da
\textbf{entrambi}.

\begin{center}
\begin{forest} classalbero
[Sistema nervoso
  [Centrale (SNC) [Encefalo][Midollo spinale]]
  [Periferico (SNP)
    [{Afferente\\(input)}
      [Sensibilità somatica][Organi di senso speciale][Sensibilità viscerale]]
    [{Efferente\\(output)}
      [Somatico [Muscolo scheletrico]]
      [{Autonomo\\(simpatico e\\parasimpatico)}
        [{Muscolo cardiaco,\\muscolo liscio,\\ghiandole}]
        [{Sistema nervoso enterico\\e tratto gastrointestinale}]]]]
]
\end{forest}
\end{center}

% ---------------------------------------------------------------------------
\section{La glia}

Nel sistema nervoso, sia centrale sia periferico, accanto ai neuroni si trovano le
\textbf{cellule della glia}, che svolgono compiti diversi.
\begin{itemize}
  \item nel \textbf{SNC}: cellule ependimali, astrociti, oligodendrociti, cellule della microglia;
  \item nel \textbf{SNP}: cellule di Schwann e cellule satelliti.
\end{itemize}

Due sono le funzioni fondamentali illustrate come esempio.
\begin{enumerate}
  \item \textbf{azione protettiva} (astrociti): si pongono tra il sistema vascolare (il capillare)
    e i neuroni, impedendo il contatto diretto tra sangue e cellule nervose $\rightarrow$ evitano
    che sostanze tossiche raggiungano i neuroni, che sono tra le cellule più delicate
    dell'organismo;
  \item \textbf{guaina mielinica} (oligodendrociti): si avvolgono attorno all'assone e formano una
    \textbf{barriera isolante}, interrotta solo a livello dei \textbf{nodi di Ranvier}, dove
    avviene la conduzione del potenziale d'azione.
\end{enumerate}

\subsection{Funzioni delle cellule gliali}

\rubrica{Sistema nervoso centrale}
\begin{itemize}
  \item \textbf{oligodendrociti}: formano le guaine mieliniche;
  \item \textbf{astrociti}:
    \begin{itemize}
      \item creano un \textbf{supporto} per il SNC, una sorta di rete su cui i neuroni poggiano;
      \item contribuiscono a formare la \textbf{barriera ematoencefalica}, frapponendosi tra
        capillare e neurone;
      \item secernono \textbf{fattori neurotrofici};
      \item \textbf{captano $K^+$ e neurotrasmettitori}: sono fondamentali nella fase finale della
        trasmissione sinaptica, perché ricaptano il potassio al termine del potenziale d'azione
        (coadiuvando la pompa $Na^+/K^+$) e soprattutto lo eliminano dallo spazio sinaptico,
        evitando che la cellula diventi ipereccitabile; ricaptano inoltre il neurotrasmettitore,
        impedendo che resti nello spazio sinaptico e si rileghi ai recettori dando origine a un
        \textit{falso messaggio} che la cellula presinaptica non intendeva più trasferire;
    \end{itemize}
  \item \textbf{microglia}: cellule di difesa specializzate, rimuovono cellule danneggiate e agenti
    estranei;
  \item \textbf{cellule ependimali}: creano barriere fra i compartimenti.
\end{itemize}

\rubrica{Sistema nervoso periferico}
\begin{itemize}
  \item \textbf{cellule satelliti}: sostengono i corpi cellulari, con un ruolo analogo a quello
    degli astrociti nel SNC;
  \item \textbf{cellule di Schwann}: formano le guaine mieliniche in periferia e secernono
    \textbf{fattori neurotrofici}. Il più importante è il \textbf{fattore di crescita nervoso}
    (\textit{nerve growth factor}, NGF), scoperta premiata con il Nobel. Le cellule di Schwann
    fanno da \textbf{guida per il cono di crescita dell'assone}, sia nella vita fetale (sviluppo)
    sia nella vita adulta (in caso di lesione dell'assone), indirizzandolo verso la cellula
    \textit{target} con cui l'accoppiamento è geneticamente determinato: così un comando che parte
    dall'$\alpha$-motoneurone giunge proprio a un determinato muscolo e non a un altro.
\end{itemize}

% ---------------------------------------------------------------------------
\section{I neuroni}

I neuroni sono di diverso tipo e di diversa \textbf{morfologia}; la morfologia ne determina la
funzione, per cui dall'aspetto di un neurone si può risalire a ciò che fa.

\rubrica{Classificazione funzionale}
\begin{enumerate}
  \item \textbf{neuroni afferenti}: portano le informazioni dalla periferia del corpo al SNC;
  \item \textbf{interneuroni}: connettono tra loro i vari neuroni;
  \item \textbf{neuroni efferenti}: danno origine a un assone che lascia il sistema nervoso per
    andare in periferia.
\end{enumerate}

\rubrica{Neuroni afferenti}
\begin{itemize}
  \item \textbf{pseudounipolare}, detto anche \textbf{a T}: tipico del sistema somatosensoriale; il
    corpo cellulare si trova nella catena dei gangli midollari e l'assone si dirama in due branche,
    una diretta al SNC (in particolare al midollo spinale, \textit{assone centrale}) e una diretta
    alla periferia a raccogliere gli stimoli sensoriali (\textit{assone periferico});
  \item \textbf{bipolare}: bipolare vero e proprio, con due prolungamenti assonici; proprio dei
    sistemi visivo e olfattivo.
\end{itemize}

\rubrica{Interneuroni}
\begin{itemize}
  \item \textbf{anassonico};
  \item \textbf{multipolare}: p.\,es.\ neuroni con \textbf{albero dendritico molto sviluppato}, il
    cui compito principale è ricevere informazioni da moltissimi altri neuroni, elaborarle e
    trasferire l'informazione finale attraverso il proprio assone (p.\,es.\ le \textbf{cellule di
    Purkinje} del cervelletto).
\end{itemize}

\rubrica{Neurone efferente}
Tipico \textbf{motoneurone}, multipolare: corpo cellulare con albero dendritico piuttosto
sviluppato, situato nelle \textbf{corna grigie del midollo spinale}; l'assone lascia il midollo
spinale ed è \textbf{mielinizzato} (guaina mielinica prodotta dalle cellule di Schwann, nodi di
Ranvier), e il suo \textbf{terminale assonico} fa sinapsi p.\,es.\ con una cellula muscolare per
indurne la contrazione.

% ---------------------------------------------------------------------------
\section{I circuiti neuronali}

Per svolgere attività così complesse i neuroni sono organizzati in \textbf{circuiti nervosi}. Ne
esistono forme elementari.

\begin{enumerate}
  \item \textbf{divergenza}: l'informazione contenuta in un neurone è trasferita a più neuroni,
    perché il suo assone forma delle \textbf{collaterali} che fanno sinapsi con neuroni differenti;
    se a loro volta questi divergono, l'informazione di partenza raggiunge un numero ancora
    maggiore di cellule $\rightarrow$ il segnale è \textbf{amplificato} e diffuso;
  \item \textbf{convergenza}: più informazioni arrivano allo stesso neurone. P.\,es.\ nelle corna
    grigie del midollo spinale un \textbf{motoneurone} riceve contemporaneamente l'afferenza dalla
    periferia sensoriale (fibra afferente della radice dorsale), quella di un \textbf{interneurone}
    e quella di una \textbf{via motoria discendente}; l'$\alpha$-motoneurone deve elaborare i tre
    tipi di informazione facendo una \textbf{sommazione spaziale} e determinando, al
    \textbf{monticolo assonico}, la comparsa o meno di un potenziale d'azione, oppure la variazione
    della \textbf{frequenza} dei potenziali d'azione diretti al muscolo effettore;
  \item \textbf{inversione del segnale}: la stessa fibra afferente eccita direttamente un neurone,
    ma con una sua collaterale fa sinapsi con un neurone \textbf{inibitorio}; a seconda del neurone
    con cui fa sinapsi, e in funzione delle informazioni discendenti ricevute, la stessa fibra fa
    uscire o non fa uscire l'informazione;
  \item \textbf{riverberazione del segnale}: il neurone eccitato, oltre a portare l'informazione ad
    altri neuroni, invia una o due collaterali a uno o due \textbf{interneuroni} che
    \textbf{riportano l'informazione a lui}. Se gli interneuroni sono eccitatori si ha
    \textbf{amplificazione} del segnale (maggiore facilità di passaggio); se sono inibitori si ha
    un \textbf{controllo} (più il neurone eccita l'interneurone, più viene inibito)
    $\rightarrow$ è possibile \textbf{modulare} quanto segnale giunge al muscolo, cioè il tono
    muscolare o la forza di contrazione in quel momento.
\end{enumerate}

Convergenza e divergenza possono coesistere sullo stesso neurone: esso riceve ed elabora una serie
di informazioni convergenti e ne distribuisce l'esito, in modo divergente, a molti altri neuroni.
A seconda di come si compongono le sinapsi eccitatorie e inibitorie del circuito si ottengono
segnali di blocco per \textbf{inibizione anterograda} o \textbf{retrograda}.

\rubrica{Ordine di grandezza}
Quelli illustrati sono esempi molto semplici: ogni neurone può ricevere \textbf{migliaia di
contatti sinaptici} nello stesso istante e altrettanti trasferirne $\rightarrow$ enorme lavoro e
consumo di energia anche per movimenti o pensieri molto semplici.

% ---------------------------------------------------------------------------
\section{Classificazione delle fibre nervose}

Le informazioni viaggiano a velocità variabile a seconda della fibra nervosa. Le fibre sono
classificate in base a \textbf{diametro} e \textbf{mielinizzazione}: il potenziale d'azione viaggia
tanto più velocemente quanto più grande è il diametro dell'assone e quanto più la fibra è
mielinizzata (la conduzione è \textbf{saltatoria}, di nodo di Ranvier in nodo di Ranvier).

\begin{table}[htbp]
  \centering
  \small
  \renewcommand{\arraystretch}{1.3}
  \begin{tabularx}{\textwidth}{|X|c|c|c|c|}
    \hline
    \textbf{Classe} & $A\alpha$ & $A\beta$ & $A\delta$ & $C$ \\
    \hline
    Diametro ($\mu$m) & 13--20 & 6--12 & 1--5 & 0,2--1,5 \\
    \hline
    Velocità di conduzione (m/s) & 78--120 & 36--72 & 6--30 & 0,5--2 \\
    \hline
  \end{tabularx}
\end{table}

\begin{itemize}
  \item $A\alpha$: sono gli \textbf{$\alpha$-motoneuroni}, i neuroni più grandi, velocissimi;
  \item $A\beta$ e $A\delta$: diametro dell'assone via via minore, ma \textbf{sempre mielinizzate},
    quindi con velocità di conduzione ancora discreta;
  \item $C$: le più lente e le più piccole di diametro; conducono p.\,es.\ la \textbf{sensazione
    dolorifica}.
\end{itemize}
